\IfFileExists{stacks-project.cls}{%
\documentclass{stacks-project}
}{%
\documentclass{amsart}
}

% For dealing with references we use the comment environment
\usepackage{verbatim}
\newenvironment{reference}{\comment}{\endcomment}
%\newenvironment{reference}{}{}
\newenvironment{slogan}{\comment}{\endcomment}
\newenvironment{history}{\comment}{\endcomment}

% For commutative diagrams we use Xy-pic
\usepackage[all]{xy}

% We use 2cell for 2-commutative diagrams.
\xyoption{2cell}
\UseAllTwocells

% We use multicol for the list of chapters between chapters
\usepackage{multicol}

% This is generally recommended for better output
\usepackage{lmodern}
\usepackage[T1]{fontenc}

% For cross-file-references
\usepackage{xr-hyper}

% Package for hypertext links:
\usepackage{hyperref}

% For any local file, say "hello.tex" you want to link to please
% use \externaldocument[hello-]{hello}
\externaldocument[introduction-]{introduction}
\externaldocument[conventions-]{conventions}
\externaldocument[sets-]{sets}
\externaldocument[categories-]{categories}
\externaldocument[topology-]{topology}
\externaldocument[sheaves-]{sheaves}
\externaldocument[sites-]{sites}
\externaldocument[stacks-]{stacks}
\externaldocument[fields-]{fields}
\externaldocument[algebra-]{algebra}
\externaldocument[brauer-]{brauer}
\externaldocument[homology-]{homology}
\externaldocument[derived-]{derived}
\externaldocument[simplicial-]{simplicial}
\externaldocument[more-algebra-]{more-algebra}
\externaldocument[smoothing-]{smoothing}
\externaldocument[modules-]{modules}
\externaldocument[sites-modules-]{sites-modules}
\externaldocument[injectives-]{injectives}
\externaldocument[cohomology-]{cohomology}
\externaldocument[sites-cohomology-]{sites-cohomology}
\externaldocument[dga-]{dga}
\externaldocument[dpa-]{dpa}
\externaldocument[sdga-]{sdga}
\externaldocument[hypercovering-]{hypercovering}
\externaldocument[schemes-]{schemes}
\externaldocument[constructions-]{constructions}
\externaldocument[properties-]{properties}
\externaldocument[morphisms-]{morphisms}
\externaldocument[coherent-]{coherent}
\externaldocument[divisors-]{divisors}
\externaldocument[limits-]{limits}
\externaldocument[varieties-]{varieties}
\externaldocument[topologies-]{topologies}
\externaldocument[descent-]{descent}
\externaldocument[perfect-]{perfect}
\externaldocument[more-morphisms-]{more-morphisms}
\externaldocument[flat-]{flat}
\externaldocument[groupoids-]{groupoids}
\externaldocument[more-groupoids-]{more-groupoids}
\externaldocument[etale-]{etale}
\externaldocument[chow-]{chow}
\externaldocument[intersection-]{intersection}
\externaldocument[pic-]{pic}
\externaldocument[weil-]{weil}
\externaldocument[adequate-]{adequate}
\externaldocument[dualizing-]{dualizing}
\externaldocument[duality-]{duality}
\externaldocument[discriminant-]{discriminant}
\externaldocument[derham-]{derham}
\externaldocument[local-cohomology-]{local-cohomology}
\externaldocument[algebraization-]{algebraization}
\externaldocument[curves-]{curves}
\externaldocument[resolve-]{resolve}
\externaldocument[models-]{models}
\externaldocument[functors-]{functors}
\externaldocument[equiv-]{equiv}
\externaldocument[pione-]{pione}
\externaldocument[etale-cohomology-]{etale-cohomology}
\externaldocument[proetale-]{proetale}
\externaldocument[relative-cycles-]{relative-cycles}
\externaldocument[more-etale-]{more-etale}
\externaldocument[trace-]{trace}
\externaldocument[crystalline-]{crystalline}
\externaldocument[spaces-]{spaces}
\externaldocument[spaces-properties-]{spaces-properties}
\externaldocument[spaces-morphisms-]{spaces-morphisms}
\externaldocument[decent-spaces-]{decent-spaces}
\externaldocument[spaces-cohomology-]{spaces-cohomology}
\externaldocument[spaces-limits-]{spaces-limits}
\externaldocument[spaces-divisors-]{spaces-divisors}
\externaldocument[spaces-over-fields-]{spaces-over-fields}
\externaldocument[spaces-topologies-]{spaces-topologies}
\externaldocument[spaces-descent-]{spaces-descent}
\externaldocument[spaces-perfect-]{spaces-perfect}
\externaldocument[spaces-more-morphisms-]{spaces-more-morphisms}
\externaldocument[spaces-flat-]{spaces-flat}
\externaldocument[spaces-groupoids-]{spaces-groupoids}
\externaldocument[spaces-more-groupoids-]{spaces-more-groupoids}
\externaldocument[bootstrap-]{bootstrap}
\externaldocument[spaces-pushouts-]{spaces-pushouts}
\externaldocument[spaces-chow-]{spaces-chow}
\externaldocument[groupoids-quotients-]{groupoids-quotients}
\externaldocument[spaces-more-cohomology-]{spaces-more-cohomology}
\externaldocument[spaces-simplicial-]{spaces-simplicial}
\externaldocument[spaces-duality-]{spaces-duality}
\externaldocument[formal-spaces-]{formal-spaces}
\externaldocument[restricted-]{restricted}
\externaldocument[spaces-resolve-]{spaces-resolve}
\externaldocument[formal-defos-]{formal-defos}
\externaldocument[defos-]{defos}
\externaldocument[cotangent-]{cotangent}
\externaldocument[examples-defos-]{examples-defos}
\externaldocument[algebraic-]{algebraic}
\externaldocument[examples-stacks-]{examples-stacks}
\externaldocument[stacks-sheaves-]{stacks-sheaves}
\externaldocument[criteria-]{criteria}
\externaldocument[artin-]{artin}
\externaldocument[quot-]{quot}
\externaldocument[stacks-properties-]{stacks-properties}
\externaldocument[stacks-morphisms-]{stacks-morphisms}
\externaldocument[stacks-limits-]{stacks-limits}
\externaldocument[stacks-cohomology-]{stacks-cohomology}
\externaldocument[stacks-perfect-]{stacks-perfect}
\externaldocument[stacks-introduction-]{stacks-introduction}
\externaldocument[stacks-more-morphisms-]{stacks-more-morphisms}
\externaldocument[stacks-geometry-]{stacks-geometry}
\externaldocument[moduli-]{moduli}
\externaldocument[moduli-curves-]{moduli-curves}
\externaldocument[examples-]{examples}
\externaldocument[exercises-]{exercises}
\externaldocument[guide-]{guide}
\externaldocument[desirables-]{desirables}
\externaldocument[coding-]{coding}
\externaldocument[obsolete-]{obsolete}
\externaldocument[fdl-]{fdl}
\externaldocument[index-]{index}

% Theorem environments.
%
\theoremstyle{plain}
\newtheorem{theorem}[subsection]{Theorem}
\newtheorem{proposition}[subsection]{Proposition}
\newtheorem{lemma}[subsection]{Lemma}

\theoremstyle{definition}
\newtheorem{definition}[subsection]{Definition}
\newtheorem{example}[subsection]{Example}
\newtheorem{exercise}[subsection]{Exercise}
\newtheorem{situation}[subsection]{Situation}

\theoremstyle{remark}
\newtheorem{remark}[subsection]{Remark}
\newtheorem{remarks}[subsection]{Remarks}

\numberwithin{equation}{subsection}

% Macros
%
\def\lim{\mathop{\mathrm{lim}}\nolimits}
\def\colim{\mathop{\mathrm{colim}}\nolimits}
\def\Spec{\mathop{\mathrm{Spec}}}
\def\Hom{\mathop{\mathrm{Hom}}\nolimits}
\def\Ext{\mathop{\mathrm{Ext}}\nolimits}
\def\SheafHom{\mathop{\mathcal{H}\!\mathit{om}}\nolimits}
\def\SheafExt{\mathop{\mathcal{E}\!\mathit{xt}}\nolimits}
\def\Sch{\mathit{Sch}}
\def\Mor{\mathop{\mathrm{Mor}}\nolimits}
\def\Ob{\mathop{\mathrm{Ob}}\nolimits}
\def\Sh{\mathop{\mathit{Sh}}\nolimits}
\def\NL{\mathop{N\!L}\nolimits}
\def\CH{\mathop{\mathrm{CH}}\nolimits}
\def\proetale{{pro\text{-}\acute{e}tale}}
\def\etale{{\acute{e}tale}}
\def\QCoh{\mathit{QCoh}}
\def\Ker{\mathop{\mathrm{Ker}}}
\def\Im{\mathop{\mathrm{Im}}}
\def\Coker{\mathop{\mathrm{Coker}}}
\def\Coim{\mathop{\mathrm{Coim}}}

% Boxtimes
%
\DeclareMathSymbol{\boxtimes}{\mathbin}{AMSa}{"02}

%
% Macros for moduli stacks/spaces
%
\def\QCohstack{\mathcal{QC}\!\mathit{oh}}
\def\Cohstack{\mathcal{C}\!\mathit{oh}}
\def\Spacesstack{\mathcal{S}\!\mathit{paces}}
\def\Quotfunctor{\mathrm{Quot}}
\def\Hilbfunctor{\mathrm{Hilb}}
\def\Curvesstack{\mathcal{C}\!\mathit{urves}}
\def\Polarizedstack{\mathcal{P}\!\mathit{olarized}}
\def\Complexesstack{\mathcal{C}\!\mathit{omplexes}}
% \Pic is the operator that assigns to X its picard group, usage \Pic(X)
% \Picardstack_{X/B} denotes the Picard stack of X over B
% \Picardfunctor_{X/B} denotes the Picard functor of X over B
\def\Pic{\mathop{\mathrm{Pic}}\nolimits}
\def\Picardstack{\mathcal{P}\!\mathit{ic}}
\def\Picardfunctor{\mathrm{Pic}}
\def\Deformationcategory{\mathcal{D}\!\mathit{ef}}
\usepackage{cleveref}
\crefname{originallemma}{original lemma}{original lemmas}
\Crefname{originallemma}{Original lemma}{Original lemmas}
\makeatletter
\newlabel{coherent-lemma-cech-cohomology-quasi-coherent}{{01XD}{source}{}{sourcepointer-1}{}}
\newlabel{coherent-lemma-cech-cohomology-quasi-coherent@cref}{{[originallemma][0][]01XD}{[1][0][]source}{}{}{}}
\makeatother
\crefname{originallemma}{original lemma}{original lemmas}
\Crefname{originallemma}{Original lemma}{Original lemmas}
\makeatletter
\newlabel{derived-lemma-morphisms-lift}{{013P}{source}{}{sourcepointer-2}{}}
\newlabel{derived-lemma-morphisms-lift@cref}{{[originallemma][0][]013P}{[1][0][]source}{}{}{}}
\makeatother
\crefname{originallemma}{original lemma}{original lemmas}
\Crefname{originallemma}{Original lemma}{Original lemmas}
\makeatletter
\newlabel{lemma-standard-covering-Cech-quasi-coherent}{{03FJ}{source}{}{sourcepointer-3}{}}
\newlabel{lemma-standard-covering-Cech-quasi-coherent@cref}{{[originallemma][0][]03FJ}{[1][0][]source}{}{}{}}
\makeatother
\crefname{originallemma}{original lemma}{original lemmas}
\Crefname{originallemma}{Original lemma}{Original lemmas}
\makeatletter
\newlabel{sites-cohomology-lemma-cech-spectral-sequence-application}{{03F7}{source}{}{sourcepointer-4}{}}
\newlabel{sites-cohomology-lemma-cech-spectral-sequence-application@cref}{{[originallemma][0][]03F7}{[1][0][]source}{}{}{}}
\makeatother
\crefname{originallemma}{original lemma}{original lemmas}
\Crefname{originallemma}{Original lemma}{Original lemmas}
\makeatletter
\newlabel{sites-cohomology-lemma-cech-vanish-collection}{{03F9}{source}{}{sourcepointer-5}{}}
\newlabel{sites-cohomology-lemma-cech-vanish-collection@cref}{{[originallemma][0][]03F9}{[1][0][]source}{}{}{}}
\makeatother
\title{Cohomology of a quasi-coherent sheaf is independent of the topology used}
\author{The Stacks Project Authors}
\date{stacks-project:descent.tex:local-snapshot-20261012; GNU FDL 1.2 or later}
\begin{document}
\maketitle
\noindent\textit{Editorial note.} All delivered mathematical text is the author's own, copied from the pinned Stacks Project chapter file descent.tex; no statement or proof step was rewritten, reordered or removed.\par\medskip
\noindent\textit{Editorial note.} The statement carries the author's own slogan; the chapter's preamble defines slogan environments as comment environments, so the slogan is part of the counted source text and is not typeset.\par\medskip
\noindent\textit{Editorial note.} The proof contains the author's sentence suggesting that the reader may skip the rest of the proof or rewrite it for a clearer exposition; the sentence is reproduced verbatim, and the derivation that follows is complete as written, so it is a reading suggestion rather than an omission.\par\medskip
\noindent\textit{Editorial note.} The counted spans carry 5 navigation stubs; no bibliography citation occurs in the spans.\par\medskip
\noindent\textit{Editorial note.} The preamble span points at the chapter's own preamble.tex shipped with the frozen source, used verbatim as staged.\par\medskip
\noindent\textit{Allowed standard prerequisites (editorial).}\begin{itemize}
\item Vanishing criterion for cohomology on a basis of objects with a cofinal system of coverings. Cohomology on Sites, Lemma sites-cohomology-lemma-cech-vanish-collection, invoked at proof lines 3018-3024. Supplies the vanishing of higher cohomology of the associated module on affines.
\item Standard affine coverings compute the Cech cohomology of a quasi-coherent module. Same chapter, Lemma lemma-standard-covering-Cech-quasi-coherent, invoked at proof lines 3024-3026. Supplies assumption (3) of the vanishing criterion.
\item Cech spectral sequence and its application to a covering. Cohomology on Sites, Lemma sites-cohomology-lemma-cech-spectral-sequence-application, invoked at proof lines 3039-3046. Supplies the comparison of Cech and ordinary cohomology on separated objects.
\item Lifting a map of complexes out of an exact complex of injectives. Derived Categories, Lemma derived-lemma-morphisms-lift, invoked at proof lines 3066-3069. Supplies the comparison map between the two derived global section complexes.
\end{itemize}\medskip
\ifdefined\hypertarget\hypertarget{sourcepointer-1}{}\fi
\noindent\textit{Source reference.} Original lemma 01XD: sources/stacks-chow/tags.txt, source lines 2514--2514. Pointer outside the counted span, invoked once for context; the load-bearing argument of the counted statement is the author's proof reproduced in full in the proof span.\par\medskip
\ifdefined\hypertarget\hypertarget{sourcepointer-2}{}\fi
\noindent\textit{Source reference.} Original lemma 013P: sources/stacks-chow/tags.txt, source lines 1444--1444. Pointer outside the counted span, invoked once for context; the load-bearing argument of the counted statement is the author's proof reproduced in full in the proof span.\par\medskip
\ifdefined\hypertarget\hypertarget{sourcepointer-3}{}\fi
\noindent\textit{Source reference.} Original lemma 03FJ: sources/stacks-descent/pinned-original-source.tex, source lines 2948--2948. Pointer outside the counted span, invoked once for context; the load-bearing argument of the counted statement is the author's proof reproduced in full in the proof span.\par\medskip
\ifdefined\hypertarget\hypertarget{sourcepointer-4}{}\fi
\noindent\textit{Source reference.} Original lemma 03F7: sources/stacks-chow/tags.txt, source lines 4482--4482. Pointer outside the counted span, invoked once for context; the load-bearing argument of the counted statement is the author's proof reproduced in full in the proof span.\par\medskip
\ifdefined\hypertarget\hypertarget{sourcepointer-5}{}\fi
\noindent\textit{Source reference.} Original lemma 03F9: sources/stacks-chow/tags.txt, source lines 4484--4484. Pointer outside the counted span, invoked once for context; the load-bearing argument of the counted statement is the author's proof reproduced in full in the proof span.\par\medskip
\begin{proposition}
\label{proposition-same-cohomology-quasi-coherent}
\begin{slogan}
Cohomology of quasi-coherent sheaves is the same no matter which
topology you use.
\end{slogan}
Let $S$ be a scheme. Let $\mathcal{F}$ be a quasi-coherent sheaf on $S$.
Let $\tau \in \{Zariski, \linebreak[0] \etale, \linebreak[0]
smooth, \linebreak[0] syntomic, \linebreak[0] fppf\}$.
\begin{enumerate}
\item There is a canonical isomorphism
$$
H^q(S, \mathcal{F}) = H^q((\Sch/S)_\tau, \mathcal{F}^a).
$$
\item There are canonical isomorphisms
$$
H^q(S, \mathcal{F}) =
H^q(S_{Zar}, \mathcal{F}^a) =
H^q(S_\etale, \mathcal{F}^a).
$$
\end{enumerate}
\end{proposition}

\begin{proof}
The result for $q = 0$ is clear from the definition of $\mathcal{F}^a$.
Let $\mathcal{C} = (\Sch/S)_\tau$, or $\mathcal{C} = S_\etale$,
or $\mathcal{C} = S_{Zar}$.

\medskip\noindent
We are going to apply
Cohomology on Sites,
Lemma \ref{sites-cohomology-lemma-cech-vanish-collection}
with $\mathcal{F} = \mathcal{F}^a$,
$\mathcal{B} \subset \Ob(\mathcal{C})$ the set of affine schemes
in $\mathcal{C}$, and $\text{Cov} \subset \text{Cov}_\mathcal{C}$ the
set of standard affine $\tau$-coverings. Assumption (3) of
the lemma is satisfied by
Lemma \ref{lemma-standard-covering-Cech-quasi-coherent}.
Hence we conclude that $H^p(U, \mathcal{F}^a) = 0$ for every
affine object $U$ of $\mathcal{C}$.

\medskip\noindent
Next, let $U \in \Ob(\mathcal{C})$ be any separated object.
Denote $f : U \to S$ the structure morphism.
Let $U = \bigcup U_i$ be an affine open covering.
We may also think of this as a $\tau$-covering
$\mathcal{U} = \{U_i \to U\}$ of $U$ in $\mathcal{C}$.
Note that
$U_{i_0} \times_U \ldots \times_U U_{i_p} =
U_{i_0} \cap \ldots \cap U_{i_p}$ is affine as we assumed $U$ separated.
By
Cohomology on Sites, Lemma
\ref{sites-cohomology-lemma-cech-spectral-sequence-application}
and the result above we see that
$$
H^p(U, \mathcal{F}^a) = \check{H}^p(\mathcal{U}, \mathcal{F}^a)
= H^p(U, f^*\mathcal{F})
$$
the last equality by
Cohomology of Schemes, Lemma
\ref{coherent-lemma-cech-cohomology-quasi-coherent}.
In particular, if $S$ is separated we can take $U = S$ and
$f = \text{id}_S$ and the proposition is proved.
We suggest the reader skip the rest of the proof (or rewrite it
to give a clearer exposition).

\medskip\noindent
Choose an injective resolution $\mathcal{F} \to \mathcal{I}^\bullet$ on $S$.
Choose an injective resolution $\mathcal{F}^a \to \mathcal{J}^\bullet$
on $\mathcal{C}$. Denote $\mathcal{J}^n|_S$ the restriction of $\mathcal{J}^n$
to opens of $S$; this is a sheaf on the topological space $S$ as open
coverings are $\tau$-coverings. We get a complex
$$
0 \to \mathcal{F} \to \mathcal{J}^0|_S \to \mathcal{J}^1|_S \to \ldots
$$
which is exact since its sections over any affine open $U \subset S$
is exact (by the vanishing of $H^p(U, \mathcal{F}^a)$, $p > 0$ seen
above). Hence by
Derived Categories, Lemma \ref{derived-lemma-morphisms-lift}
there exists map of complexes
$\mathcal{J}^\bullet|_S \to \mathcal{I}^\bullet$ which in particular
induces a map
$$
R\Gamma(\mathcal{C}, \mathcal{F}^a)
=
\Gamma(S, \mathcal{J}^\bullet)
\longrightarrow
\Gamma(S, \mathcal{I}^\bullet)
=
R\Gamma(S, \mathcal{F}).
$$
Taking cohomology gives the map
$H^n(\mathcal{C}, \mathcal{F}^a) \to H^n(S, \mathcal{F})$ which
we have to prove is an isomorphism.
Let $\mathcal{U} : S = \bigcup U_i$ be an affine open covering
which we may think of as a $\tau$-covering also.
By the above we get a map of double complexes
$$
\check{\mathcal{C}}^\bullet(\mathcal{U}, \mathcal{J})
=
\check{\mathcal{C}}^\bullet(\mathcal{U}, \mathcal{J}|_S)
\longrightarrow
\check{\mathcal{C}}^\bullet(\mathcal{U}, \mathcal{I}).
$$
This map induces a map of spectral sequences
$$
{}^\tau\! E_2^{p, q} = \check{H}^p(\mathcal{U}, \underline{H}^q(\mathcal{F}^a))
\longrightarrow
E_2^{p, q} = \check{H}^p(\mathcal{U}, \underline{H}^q(\mathcal{F}))
$$
The first spectral sequence converges to
$H^{p + q}(\mathcal{C}, \mathcal{F})$ and the second to
$H^{p + q}(S, \mathcal{F})$. On the other hand, we have seen
that the induced maps ${}^\tau\! E_2^{p, q} \to E_2^{p, q}$ are
bijections (as all the intersections are separated being opens in affines).
Whence also the maps $H^n(\mathcal{C}, \mathcal{F}^a) \to H^n(S, \mathcal{F})$
are isomorphisms, and we win.
\end{proof}

\end{document}
